% !TEX program = xelatex
\documentclass[14pt]{extarticle}

\usepackage[
  a4paper,
  left=2.5cm,
  right=1.5cm,
  top=2cm,
  bottom=2cm
]{geometry}

\usepackage[ukrainian]{babel}
\usepackage{fontspec}
\setmainfont{Times New Roman}

\usepackage{setspace}
\onehalfspacing
\setlength{\parindent}{1.27cm}
\setlength{\parskip}{0pt}
\usepackage{microtype}

\usepackage{titlesec}
\titleformat{\section}
  {\bfseries\fontsize{16pt}{24pt}\selectfont\centering}
  {\thesection}{0.5em}{}
\titlespacing*{\section}{0pt}{20pt}{6pt}

\titleformat{\subsection}
  {\bfseries\fontsize{14pt}{21pt}\selectfont}
  {\thesubsection}{0.5em}{}
\titlespacing*{\subsection}{1.27cm}{6pt}{6pt}

\usepackage{enumitem}
\setlist[itemize,1]{
  label=$\bullet$,
  leftmargin=1.27cm,
  labelwidth=0.5cm,
  labelsep=0.3cm,
  itemindent=0pt,
  itemsep=3pt,
  topsep=6pt,
  parsep=0pt,
  partopsep=0pt
}
\setlist[enumerate,1]{
  label=\arabic*.,
  leftmargin=2.54cm,
  labelwidth=0.63cm,
  labelsep=0.54cm,
  itemindent=0pt,
  itemsep=0pt,
  topsep=0pt,
  parsep=0pt,
  partopsep=0pt
}

\usepackage{array}
\usepackage{longtable}
\usepackage{booktabs}
\usepackage{tabularx}
\usepackage{multirow}
\usepackage{graphicx}
\usepackage[table]{xcolor}

\usepackage{tikz}
\usetikzlibrary{arrows.meta,positioning,calc}
\usepackage{pgfgantt}

\usepackage[hidelinks,breaklinks=true]{hyperref}

\begin{document}
\sloppy
\emergencystretch=3em
\setlength{\tabcolsep}{4pt}

% ==================== ТИТУЛ ====================
\begin{center}
\textbf{Лабораторна робота №5}

\medskip

\textbf{\MakeUppercase{Планування та побудова «Дорожньої карти» проєкту. Побудова матриці відповідальності та комунікаційної матриці проєкту. Побудова діаграми Ганта.}}
\end{center}

\medskip

\textbf{Мета роботи:} навчитися планувати проєкт, будувати мережеву діаграму, знаходити критичний шлях, формувати матрицю комунікацій та діаграму Ганта.

\textbf{Проєкт:} <<Система для автоматичної нормалізації суржику та діалектних форм української мови>>.

% ==================== РОЗДІЛ 1: КОНТРОЛЬНИЙ СПИСОК ====================
\section*{1. Контрольний список планування проєкту}

Контрольний список охоплює 16 обов'язкових кроків планування відповідно до методології PMBOK (Project Management Body of Knowledge).

\vspace{6pt}
\noindent\textbf{Таблиця 1 --- Контрольний список планування проєкту}

\begin{longtable}{|>{\centering\arraybackslash}p{0.05\textwidth}|>{\raggedright\arraybackslash}p{0.42\textwidth}|>{\raggedright\arraybackslash}p{0.14\textwidth}|>{\raggedright\arraybackslash}p{0.24\textwidth}|}
\hline
\textbf{№} & \textbf{Крок планування} & \textbf{Статус} & \textbf{Результат} \\
\hline
\endfirsthead
\hline
\textbf{№} & \textbf{Крок планування} & \textbf{Статус} & \textbf{Результат} \\
\hline
\endhead
1 & Визначення цілей і критеріїв успіху проєкту & Виконано & Цілі сформульовано у ЛР~№1 \\
\hline
2 & Ідентифікація стейкхолдерів & Виконано & Матриця 10 стейкхолдерів (ЛР~№3) \\
\hline
3 & Визначення обмежень та припущень & Виконано & Часові рамки: лютий--червень 2026 \\
\hline
4 & Розробка WBS & Виконано & Деревоподібна структура (розд.~2) \\
\hline
5 & Визначення переліку операцій & Виконано & 6 пакетів робіт (A--F) \\
\hline
6 & Визначення залежностей між операціями & Виконано & Послідовний та паралельний потік \\
\hline
7 & Оцінка тривалості операцій & Виконано & Від 2 до 6 тижнів \\
\hline
8 & Побудова мережевої діаграми & Виконано & Метод критичного шляху (розд.~3) \\
\hline
9 & Розрахунок критичного шляху & Виконано & $A\to B\to D\to E\to F$ = 19 тижнів \\
\hline
10 & Розробка матриці відповідальності (RAM) & Виконано & Таблиця RAM у ЛР~№3 \\
\hline
11 & Планування комунікацій & Виконано & Матриця комунікацій (розд.~5) \\
\hline
12 & Ідентифікація та оцінка ризиків & Виконано & Реєстр ризиків у ЛР~№4 \\
\hline
13 & Планування ресурсів & Виконано & Виконавець + науковий керівник \\
\hline
14 & Розробка бюджету & Виконано & 0~грн (академічний проєкт) \\
\hline
15 & Побудова діаграми Ганта & Виконано & Розклад проєкту (розд.~6) \\
\hline
16 & Затвердження плану та kick-off зустріч & Виконано & Погоджено з науковим керівником \\
\hline
\end{longtable}

% ==================== РОЗДІЛ 2: WBS ====================
\section*{2. Структура розбиття робіт (WBS)}

\noindent\textbf{0.\ Система автоматичної нормалізації суржику та діалектних форм}

\vspace{4pt}
\begin{enumerate}[label=\textbf{\arabic*.}, leftmargin=1.5cm, itemsep=3pt, topsep=4pt, parsep=0pt]
  \item \textbf{Управління проєктом}
  \begin{enumerate}[label=1.\arabic*., leftmargin=1.2cm, itemsep=1pt, topsep=2pt, parsep=0pt, partopsep=0pt]
    \item Координація та звітність
    \item Управління ризиками
  \end{enumerate}

  \item \textbf{Дослідження та аналіз}
  \begin{enumerate}[label=2.\arabic*., leftmargin=1.2cm, itemsep=1pt, topsep=2pt, parsep=0pt, partopsep=0pt]
    \item Огляд академічних джерел та NLP-корпусів
    \item Аналіз наявних підходів до нормалізації
  \end{enumerate}

  \item \textbf{Підготовка даних}
  \begin{enumerate}[label=3.\arabic*., leftmargin=1.2cm, itemsep=1pt, topsep=2pt, parsep=0pt, partopsep=0pt]
    \item Збір суржик-текстів та діалектних зразків
    \item Розробка правил нормалізації
    \item Анотація та верифікація датасету
  \end{enumerate}

  \item \textbf{Розробка моделей}
  \begin{enumerate}[label=4.\arabic*., leftmargin=1.2cm, itemsep=1pt, topsep=2pt, parsep=0pt, partopsep=0pt]
    \item Seq2seq базова модель
    \item Fine-tuning Lapa LLM
  \end{enumerate}

  \item \textbf{Оцінка та інтеграція}
  \begin{enumerate}[label=5.\arabic*., leftmargin=1.2cm, itemsep=1pt, topsep=2pt, parsep=0pt, partopsep=0pt]
    \item Автоматична оцінка (BLEU, chrF, WER)
    \item Ручна оцінка носіями мови
    \item Інтеграція REST API
  \end{enumerate}

  \item \textbf{Документація та захист}
  \begin{enumerate}[label=6.\arabic*., leftmargin=1.2cm, itemsep=1pt, topsep=2pt, parsep=0pt, partopsep=0pt]
    \item Дипломна робота
    \item Технічна документація
    \item Захист проєкту
  \end{enumerate}
\end{enumerate}

% ==================== РОЗДІЛ 3: МЕРЕЖЕВА ДІАГРАМА ====================
\section*{3. Мережева діаграма проєкту}

Мережева діаграма побудована методом критичного шляху (CPM). Таблиця 2 містить перелік операцій з оцінкою тривалості та залежностями.

\vspace{6pt}
\noindent\textbf{Таблиця 2 --- Перелік операцій проєкту}

\begin{longtable}{|>{\centering\arraybackslash}p{0.04\textwidth}|>{\raggedright\arraybackslash}p{0.35\textwidth}|>{\centering\arraybackslash}p{0.10\textwidth}|>{\centering\arraybackslash}p{0.13\textwidth}|>{\raggedright\arraybackslash}p{0.22\textwidth}|}
\hline
\textbf{ID} & \textbf{Операція} & \textbf{Три\-вал.,\newline тижні} & \textbf{Залежить від} & \textbf{Відповідальний} \\
\hline
\endfirsthead
\hline
\textbf{ID} & \textbf{Операція} & \textbf{Три\-вал.,\newline тижні} & \textbf{Залежить від} & \textbf{Відповідальний} \\
\hline
\endhead
A & Аналіз літератури та предметної галузі & 3 & --- & Коваль Д.П. \\
\hline
B & Підготовка та анотація датасету & 6 & A & Коваль Д.П. \\
\hline
C & Seq2seq базова модель & 3 & B & Коваль Д.П. \\
\hline
D & Fine-tuning Lapa LLM & 4 & B & Коваль Д.П. \\
\hline
E & Оцінка та інтеграція & 2 & C, D & Коваль Д.П., Сивоконь О.О. \\
\hline
F & Написання дипломної роботи & 4 & E & Коваль Д.П., Сивоконь О.О. \\
\hline
\end{longtable}

\vspace{8pt}
\noindent\textbf{Рисунок 1 --- Мережева діаграма проєкту (метод CPM)}

\vspace{8pt}
\begin{center}
\resizebox{0.98\textwidth}{!}{%
\begin{tikzpicture}[
  taskbox/.style={
    draw=gray!70, rectangle, rounded corners=3pt,
    minimum width=3.0cm, minimum height=2.8cm,
    align=center, thick
  },
  critbox/.style={
    draw=red!80!black, rectangle, rounded corners=3pt,
    minimum width=3.0cm, minimum height=2.8cm,
    align=center, thick, line width=2pt,
    fill=red!8
  },
  arr/.style={-{Stealth[length=9pt]}, thick, gray!60},
  critarr/.style={-{Stealth[length=9pt]}, thick, red!80!black, line width=2pt}
]

\node[critbox] (A) at (0, 0) {
  \large\textbf{A}\\[3pt]
  \normalsize Аналіз\\[-1pt]
  \normalsize літератури\\[5pt]
  \small ES=0 \quad EF=3\\
  \small LS=0 \quad LF=3\\[2pt]
  \small TF=0
};

\node[critbox] (B) at (5, 0) {
  \large\textbf{B}\\[3pt]
  \normalsize Підготовка\\[-1pt]
  \normalsize даних\\[5pt]
  \small ES=3 \quad EF=9\\
  \small LS=3 \quad LF=9\\[2pt]
  \small TF=0
};

\node[taskbox] (C) at (10, 2.8) {
  \large\textbf{C}\\[3pt]
  \normalsize Seq2seq\\[-1pt]
  \normalsize модель\\[5pt]
  \small ES=9 \quad EF=12\\
  \small LS=10 \; LF=13\\[2pt]
  \small TF=1
};

\node[critbox] (D) at (10, -2.8) {
  \large\textbf{D}\\[3pt]
  \normalsize Fine-tuning\\[-1pt]
  \normalsize Lapa LLM\\[5pt]
  \small ES=9 \quad EF=13\\
  \small LS=9 \quad LF=13\\[2pt]
  \small TF=0
};

\node[critbox] (E) at (15.5, 0) {
  \large\textbf{E}\\[3pt]
  \normalsize Оцінка та\\[-1pt]
  \normalsize інтеграція\\[5pt]
  \small ES=13 \; EF=15\\
  \small LS=13 \; LF=15\\[2pt]
  \small TF=0
};

\node[critbox] (F) at (20.5, 0) {
  \large\textbf{F}\\[3pt]
  \normalsize Дипломна\\[-1pt]
  \normalsize робота\\[5pt]
  \small ES=15 \; EF=19\\
  \small LS=15 \; LF=19\\[2pt]
  \small TF=0
};

% Arrows
\draw[critarr] (A.east) -- (B.west);
\draw[arr] (B.east) -- (C.west);
\draw[critarr] (B.east) -- (D.west);
\draw[arr] (C.east) -- (E.west);
\draw[critarr] (D.east) -- (E.west);
\draw[critarr] (E.east) -- (F.west);

\end{tikzpicture}%
}
\end{center}

\vspace{4pt}
\noindent{\small\textit{Умовні позначення: червона рамка --- критичні операції (TF=0); сіра рамка --- некритичні (TF$>$0).}}

% ==================== РОЗДІЛ 4: РУХ ВПЕРЕД / ЗВОРОТНІЙ ====================
\section*{4. «Рух вперед» та «Зворотній рух»}

Метод критичного шляху передбачає два проходи: \textbf{«рух вперед»} (розрахунок раннього початку ES та раннього закінчення EF) та \textbf{«зворотній рух»} (розрахунок пізнього початку LS та пізнього закінчення LF).

\vspace{6pt}
\noindent\textbf{Таблиця 3 --- Розрахунок раннього та пізнього початку і закінчення}

\begin{longtable}{|>{\centering\arraybackslash}p{0.04\textwidth}|>{\centering\arraybackslash}p{0.09\textwidth}|>{\centering\arraybackslash}p{0.06\textwidth}|>{\centering\arraybackslash}p{0.06\textwidth}|>{\centering\arraybackslash}p{0.06\textwidth}|>{\centering\arraybackslash}p{0.06\textwidth}|>{\centering\arraybackslash}p{0.06\textwidth}|>{\centering\arraybackslash}p{0.10\textwidth}|>{\centering\arraybackslash}p{0.06\textwidth}|}
\hline
\textbf{ID} & \textbf{Три\-вал.} & \textbf{ES} & \textbf{EF} & \textbf{LS} & \textbf{LF} & \textbf{TF} & \textbf{Крит.?} & \textbf{FF} \\
\hline
\endfirsthead
\hline
\textbf{ID} & \textbf{Три\-вал.} & \textbf{ES} & \textbf{EF} & \textbf{LS} & \textbf{LF} & \textbf{TF} & \textbf{Крит.?} & \textbf{FF} \\
\hline
\endhead
A & 3 & 0 & 3 & 0 & 3 & 0 & Так & 0 \\
\hline
B & 6 & 3 & 9 & 3 & 9 & 0 & Так & 0 \\
\hline
C & 3 & 9 & 12 & 10 & 13 & 1 & Ні & 1 \\
\hline
D & 4 & 9 & 13 & 9 & 13 & 0 & Так & 0 \\
\hline
E & 2 & 13 & 15 & 13 & 15 & 0 & Так & 0 \\
\hline
F & 4 & 15 & 19 & 15 & 19 & 0 & Так & 0 \\
\hline
\end{longtable}

\noindent де: \textbf{ES} --- ранній початок; \textbf{EF} --- раннє закінчення; \textbf{LS} --- пізній початок; \textbf{LF} --- пізнє закінчення; \textbf{TF} --- загальний резерв часу; \textbf{FF} --- вільний резерв часу.

\medskip

\noindent Загальна тривалість проєкту: \textbf{19 тижнів}. Критичний шлях: $\mathbf{A \to B \to D \to E \to F}$.

% ==================== РОЗДІЛ 5: МАТРИЦЯ КОМУНІКАЦІЙ ====================
\section*{5. Матриця комунікацій проєкту}

Матриця комунікацій визначає, які повідомлення, яким стейкхолдерам, у якій формі та з якою частотою передаються впродовж проєкту.

\vspace{6pt}
\noindent\textbf{Таблиця 4 --- Матриця комунікацій}

\begin{longtable}{|>{\raggedright\arraybackslash}p{0.15\textwidth}|>{\raggedright\arraybackslash}p{0.19\textwidth}|>{\raggedright\arraybackslash}p{0.14\textwidth}|>{\raggedright\arraybackslash}p{0.14\textwidth}|>{\raggedright\arraybackslash}p{0.13\textwidth}|>{\raggedright\arraybackslash}p{0.13\textwidth}|}
\hline
\textbf{Аудиторія} & \textbf{По\-ві\-дом\-лення} & \textbf{Тип} & \textbf{Частота} & \textbf{Кому\-нікатор} & \textbf{Від\-по\-ві\-даль\-ність} \\
\hline
\endfirsthead
\hline
\textbf{Аудиторія} & \textbf{По\-ві\-дом\-лення} & \textbf{Тип} & \textbf{Частота} & \textbf{Кому\-нікатор} & \textbf{Від\-по\-ві\-даль\-ність} \\
\hline
\endhead
Науковий керівник (S2) & Звіти про прогрес та технічні питання & Зустріч / email & Що\-тижнево & Коваль Д.П. & Звітність та схвалення \\
\hline
Кафедра СШІ (S3) & Проміжні результати, відповідність вимогам & Пре\-зен\-тація & Що\-місяця & Коваль Д.П. & Від\-по\-від\-ність вимогам \\
\hline
NLP-спільнота (S7) & Публікація моделей та датасетів & GitHub / звіт & За потреби & Коваль Д.П. & Від\-критий доступ \\
\hline
Носії діалектів (S5) & Збір зразків та зворотній зв'язок & Анкета & 1 раз & Коваль Д.П. & Якість даних \\
\hline
Команда Lapa LLM (S4) & Технічні консультації та доступ до моделі & Email / API & За потреби & Коваль Д.П. & Інтеграція LLM \\
\hline
EdTech-платформи (S8) & Демонстрація готової системи & Вебінар & 1 раз & Коваль Д.П.,\newline Сивоконь О.О. & Попу\-ляри\-зація \\
\hline
\end{longtable}

% ==================== РОЗДІЛ 6: ДІАГРАМА ГАНТА ====================
\section*{6. Діаграма Ганта}

Діаграма Ганта відображає часовий розклад виконання операцій проєкту. Критичні операції (TF=0) виділено червоним кольором.

\vspace{8pt}
\noindent\textbf{Рисунок 2 --- Діаграма Ганта проєкту (тижні 1--19)}

\vspace{8pt}
\begin{center}
\resizebox{\textwidth}{!}{%
\begin{ganttchart}[
  hgrid,
  vgrid={*{1}{draw=gray!30}},
  x unit=0.70cm,
  y unit title=0.80cm,
  y unit chart=0.90cm,
  title/.style={fill=gray!25, draw=black, font=\normalsize},
  title label font=\normalsize,
  bar label font=\normalsize,
  bar/.style={fill=blue!55, draw=blue!75!black},
  bar height=0.55,
  group right shift=0,
  group top shift=0.7,
  group height=0.3
]{1}{19}
  \gantttitle{Тижні проєкту}{19} \\
  \gantttitlelist{1,...,19}{1} \\
  \ganttbar[bar/.style={fill=red!60, draw=red!80!black}]
    {A: Аналіз літератури}{1}{3} \\
  \ganttbar[bar/.style={fill=red!60, draw=red!80!black}]
    {B: Підготовка даних}{4}{9} \\
  \ganttbar[bar/.style={fill=blue!55, draw=blue!75!black}]
    {C: Seq2seq модель}{10}{12} \\
  \ganttbar[bar/.style={fill=red!60, draw=red!80!black}]
    {D: Fine-tuning Lapa LLM}{10}{13} \\
  \ganttbar[bar/.style={fill=red!60, draw=red!80!black}]
    {E: Оцінка та інтеграція}{14}{15} \\
  \ganttbar[bar/.style={fill=red!60, draw=red!80!black}]
    {F: Дипломна робота}{16}{19} \\
\end{ganttchart}%
}
\end{center}

\vspace{4pt}
\noindent{\small\textit{Умовні позначення: \textcolor{red!70!black}{\rule{1em}{0.8em}} --- критичні операції ($A\to B\to D\to E\to F$); \textcolor{blue!70!black}{\rule{1em}{0.8em}} --- некритичні операції (TF$>$0).}}

\end{document}
